\cotefile{37}
\scanpage{47}
\diagram{\begin{tikzcd}[ampersand replacement=\&,column sep=large]
 E' \arrow[r,shift right=.55ex,"\varphi'"'] \& E'' \arrow[l,shift right=.55ex,"\varphi''"']
\end{tikzcd}
\quad
$\left\{\begin{aligned}
\deleted{\lambda'\colon \operatorname{id}_{E'}\longrightarrow\varphi'\varphi''}\\
\deleted{\lambda''\colon \operatorname{id}_{E''}\longrightarrow\varphi''\varphi'}
\end{aligned}\right.$}

% Several long diagonal strokes cancel this first large scheme.
\diagram{\resizebox{\linewidth}{!}{\begin{tikzpicture}
\node[anchor=south west,inner sep=0pt,outer sep=0pt] (grid) at (0,0) {\begin{tikzcd}[ampersand replacement=\&,row sep=large,column sep=huge]
 \deleted{\text{Esp fin}} \arrow[r,leftrightarrow] \arrow[d,leftrightarrow,"\deleted{\text{topologie opposée}}"'] \&
 \deleted{\text{ord fin}} \arrow[d,leftrightarrow,"\deleted{\text{ordre opposé}}"]\\
 \deleted{\text{Espaces \unsure{topologiques|top. quelconques|top. généraux} finis}} \arrow[d]
 \arrow[r,leftrightarrow,"{\deleted{X\mapsto X\ \text{avec rel. spéc}}}","{\deleted{\substack{\text{topologie}\\\text{des cribles}}\mathrel{\reflectbox{$\mapsto$}} X}}"'] \&
 \deleted{\text{ens préordonnés finis}} \arrow[d]\\
 \deleted{\text{espaces sobres finis}} \arrow[d,leftrightarrow] \&
 \deleted{\text{ens ordonnés finis}} \arrow[l,phantom,"{\leftarrow}" description]\\
 \deleted{\text{\unsure{ens. ordonné fini modulaire|ens. ordonnés finis modulaires|ens. ordonnés finis ou dual}}} \&
\end{tikzcd}};
\begin{scope}[x={(grid.south east)},y={(grid.north west)}]
 \draw[thin] (.23,1.01) -- (.05,.26) -- (.06,.28);
 \draw[thin] (.38,1.01) -- (.19,.035) -- (.205,.064);
 \draw[thin] (.53,.96) -- (.335,.005) -- (.35,.017);
 \draw[thin] (.73,.89) -- (.565,.10) -- (.577,.12);
 \draw[thin] (.91,.79) -- (.77,.27);
 \draw[thin] (.993,.63) -- (.93,.39);
\end{scope}
\end{tikzpicture}}}

\diagram{\resizebox{\linewidth}{!}{\begin{tikzcd}[ampersand replacement=\&,column sep=large,row sep=huge]
 \text{Esp fin} \arrow[rr,leftrightarrow] \arrow[d] \& \&
 \text{préord fin} \arrow[d] \& \deleted{\unsure{y\leftarrow x|y\to x}}\\
 X=(\unsure{\mathscr X|X},T)\in\text{Esp fin sob} \arrow[rr,leftrightarrow]
 \arrow[dr,leftrightarrow,"{X\to\operatorname{Ouv}(X)}","{\deleted{\text{\unsure{ens des ouverts|ens des ouverts.|ens. des ouverts}}}}"']
 \arrow[ddr,bend right=18,"{\operatorname{Top}(\mathscr X,T)}"'] \& \&
 \text{ord fin}\ni(\mathscr X,\omega)=I \arrow[dl,leftrightarrow,"\text{\unsure{ens des cri|ens des cri.|ens. des cribles}}"]
 \arrow[ddl,bend left=28,"{(X,\omega)^{\wedge}}"] \&
 \substack{x\leftarrow y\\[-1pt]\text{\unsure{spé.|spéc.}}\qquad\text{gén}}\\
 \&\text{Lat fin mod }\unsure{\sigma|\mathcal O|O}
 \arrow[d,shift left=.6ex,"{\sigma\to\sigma^{\sim}}"] \&
 \underset{\longrightarrow}{\phantom{\sigma\to\sigma^\sim}}\ \text{\unsure{top can|top canon|top coh}} \&\\
 \& \text{Top fin \unsure{spec|spéc.}}
 \arrow[u,shift left=.6ex,"{E\mapsto\operatorname{Ouv}(E)}"]
 \arrow[uul,bend left=52,"{E\to\operatorname{Point}(E)}"']
 \arrow[uurr,bend right=36,"{\underline{\underline E}\to\text{\uline{Point}}(E)}"']
 \& \&
\end{tikzcd}}}

\margin{\deleted{\unsure{ensemble polye.|ensemble polyéd.|ensemble polyèdre}}\ schéma simplicial \inserted{fini} ($=$ ens avec
\deleted{\unsure{anti|enti|ensemble}} \unsure{filtre de|filtre des|filtres des|filtrant de} parties.\\
$\deleted{x\leq y\qquad x\longrightarrow y}$}

\diagram{\fbox{\resizebox{.93\linewidth}{!}{\begin{tikzcd}[ampersand replacement=\&,row sep=large,column sep=large]
 \text{\uline{Dictionnaire}} \& \& \& \&\\[-1.5ex]
 \& X \arrow[rr,shift left=.7ex,"{\substack{\text{\uline{Point}}(X)\ \text{ou}\\\text{relation de génér.}}}"]
 \arrow[dr,shift left=.7ex,"{\operatorname{Ouv}(X)}"]
 \arrow[dddr,bend right=48,"{\operatorname{Top}(X)}"']
 \arrow[ddl,"\substack{\text{topologie}\\\text{opposée}}"'] \& \&
 I \arrow[ll,shift left=.7ex,"\text{topologie des cribles}"']
 \arrow[dl,shift right=.7ex,"{\operatorname{Crib}(I)}"']
 \arrow[ddddll,"\substack{\text{ens ord.}\\\text{opposé}}"'] \arrow[dddl,bend left=48,"{I^{\wedge}}"] \&\\
 \& \& \sigma \arrow[ul,shift left=.7ex,very thick,"{\operatorname{Fib}(\sigma)^{o}}"]
 \arrow[ur,shift right=.7ex,very thick,"{\text{\uline{Fib}}(\sigma)^{o}}"']
 \arrow[dd,shift left=.55ex,"{\sigma^{\sim}}"]
 \arrow[dddll,"\substack{\text{ens. ordonné}\\\text{opposé}}"'] \& \&\\
 {} \& \& \& \&\\
 \& \& E \arrow[uu,shift left=.55ex,"{\operatorname{Ouv}(E)}"]
 \arrow[uuul,bend left=59,"{\operatorname{Point}(E)}"']
 \arrow[uuur,bend right=64,"{\text{\uline{Point}}(E)}"']
 \arrow[ddl,very thick,bend right=10,"{\tikz[baseline=(a.base)]{\node[inner sep=0pt](a){$?E$};\draw (a.center) ellipse (.86em and .8em);}\mapsto E^{\tikz[baseline=(b.base)]{\node[inner sep=0pt](b){$\scriptstyle\mathrm{opp}?$};\draw (b.center) ellipse (1em and .6em);}}}"] \& \&\\
 {} \& {} \& \& \&\\
 \& {} \& \& \&
\end{tikzcd}}}}

\[
\begin{aligned}
 & I^{\wedge}\qquad I^{\vee}\\
\deleted{\unsure{x\leftrightarrow y|x\to y}}\quad&
 \boxed{C^{\wedge}\ \text{et}\ C^{\vee}=\unsure{(C^{o})^{\wedge}|C^{o\wedge}|C^{\circ\wedge}}}\\
 &\tikz[baseline=(q.base)]{\node[inner sep=0pt](q){\uline{topos opposé} ?};\draw (q.center) ellipse (3.85em and .9em);}
\end{aligned}
\]
% The last question is circled; the question marks by E and by opp are separately circled.
% archive sheet number: 46
\scanpage{48}
\mspage{\handwritten{19}}
% The typescript is upside down in the scan. A long pen stroke cancels the sheet.
\deleted{9.1.13. \uline{Questions ouvertes}. Soit $E$ un \uline{U}-topos.}

\deleted{a) Dans l'ensemble \inserted{ordonné} des sous-topos de\unsure{,} $E$, le Inf de deux éléments est-il distributif par rapport aux Sup finis ? Pour des résultats dans c e \deleted{\unsure{se|sens}} sens, voir 9.1.11 e) et 9.1.12 d) .}

\deleted{b) Soit $E'$ un sous-topos de $E$, distinct de \deleted{sous-topos ``vide''} $E$. Existe-t-il un sous-topos $E''$ de $E$, tel que $E''\handwritten{\cap}E'=\operatorname{Top}(\varnothing)$, $E''\handwritten{\ne}\operatorname{Top}(\varnothing)$ ?}

\deleted{\deleted{\unsure{b)|p}}c) Soient $E'$ et $E''$ deux sous-topos de $E$ tels que $E'\handwritten{\cap}E''=\operatorname{Top}(\varnothing)$. Les conditions suivantes sont-elles équivalentes: (i) $E$ est le sup de $E'$ et $E''$ dans l'ensemble des sous-topos de $E$ (on se retient à grand peine d'écrire cette condition $E=E'\handwritten{\deleted{\unsure{\cap|\vee}}\cup}E''$ ); (ii) $E''$ est \deleted{max} la plus grand parmi les sous-topos $E''_{1}$ de $E$ tels que $\operatorname{Top}(E')\handwritten{\cap}\operatorname{Top}(E''_{1})=\operatorname{Top}(\varnothing)$; (iii) $E''$ est maximal parmi les $E''_{1}$ précédents.}

\deleted{Il est \deleted{\unsure{a bit|abit|abit.}} clair que (ii)$\handwritten{\Longrightarrow}$(iii), que si a) est vérifié \unsure{en|on} a (i)$\handwritten{\Longrightarrow}$(ii), \deleted{\unsure{et que si|e que si}} \inserted{que} b) est \unsure{vérifié|vérifi} \deleted{\unsure{e|é}} \inserted{si et seulement si \deleted{on a}} on a (iii)$\handwritten{\Longrightarrow}$(i); donc une solution affirmative à a) et b) en impliquerait une pour c). Noter aussi que la réponse aux questions posées \inserted{( y compris d) qui suit)} est affirmative si $E$ est de la forme $\operatorname{Top}(X)$, $X$ un espace topologique \inserted{localement} noethérien (9.1.11 b)), \deleted{ainsi que pour \unsure{$E=G$|$E\in G$}} \deleted{ou si $E$ est} un topos fini (9.1.12).}

\deleted{d) Si $E$ est le sup des deux sous-topos $E'$ et $E''$, alors tout point de $E$ est-il isomorphe à un point provenant de $E'$ ou de $E''$ ?}

\deleted{Si la réponse à d) était affirmative, alors on aurait une réponse affirmative à a) \inserted{$(E'\ (E''\ E''')=E'\ E''\quad E'\ E''')$} pour le\deleted{s} cas de sous-topos $E',E'',E'''$ de $E$ ayant suffisamment de points (donc a) serait vrai pour $E$ \inserted{localement} noethérien, en vertu de 9.1.11 b)). Inversement, si a)}
% archive sheet number: 47
\scanpage{49}
$I$ ens ordonné

On va \unsure{interpréter|interprêter} $I^{\wedge}$ comme un $\operatorname{Top}(X)$

On a $I\simeq\operatorname{Top}(X)$ ; considérons la topologie canonique, Pour ceci, quels sont les cribles de $I$ : parties $\unsure{\sigma|\mathcal O|O}$ de $I$ telles que pour $\forall\ x\in\sigma$, les $y\leq x$ sont dans $\sigma$. On met sur $I$ la topologie définie par cet ens. de \deleted{cribles} parties (\deleted{\unsure{anti-|anti}}\unsure{appeler|appelle|pellicules} la topologie des cribles) On trouve que \deleted{\unsure{$\operatorname{Top}($|$\operatorname{Top}(I$}} l'on.
\[
\begin{aligned}
\operatorname{Ouv}(I)&\overset{\mathrm{déf}}{=}\operatorname{Crib}(I)\simeq\operatorname{Ouv}(I^{\wedge}) \qquad\text{donc}\\
{\unsure{I|\mathbb I}}^{\wedge}&\simeq\operatorname{Ouv}(I)^{\sim}=\operatorname{Top}(I).
\end{aligned}
\]

\uline{NB} La relation $y\leq x$ signifie que $y$ est une \uline{générisation} de $x$ (\unsure{non|ou|oui|$x$} : une spécialisation)

\diagram{\begin{tikzcd}[ampersand replacement=\&,row sep=large,nodes={anchor=west},column sep=0pt]
 \bullet\quad\parbox{.82\linewidth}{$E$ a suffisamment de \unsure{pts essentiels|pts essentiels,} et est engendré par ses ouverts} \arrow[d,leftrightarrow,start anchor=south west,end anchor=north west]\\
 \parbox{.86\linewidth}{$E$ est engendré par les \deleted{\unsure{objets|objet}} sous-objets de l'objet final qui sont connexes, non vides et projectifs.} \arrow[d,Leftrightarrow,start anchor=south west,end anchor=north west]\\
 \parbox{.86\linewidth}{(\deleted{\unsure{où|Que}}\unsure{si|où|Ex.} $E=C^{\wedge}$) Tous les morphismes de $C$ sont des monomorphismes} \arrow[d,Leftrightarrow,start anchor=south west,end anchor=north west]\\
 \parbox{.86\linewidth}{(si $E=\operatorname{Top}(X)$, $X$ sobre) Pour $\forall\ x\in X$, l'ens des générisations de $x$ est un voisinage de $x$.}
\end{tikzcd}}
% archive sheet number: 48
\scanpage{50}
\mspage{\handwritten{6}}
% The typescript is upside down in the scan and crossed by long diagonal pen strokes.
\deleted{Pour la première assertion, on se ramène encore, par descente et en utilisant \deleted{1}\handwritten{4}.\handwritten{3}, au cas où $k$ est séparablement clos. Procédant comme dans la démonstration de \deleted{1}\handwritten{4}.6 , on déduit de \deleted{1}\handwritten{4}.5 et de l'hypothèse sur $L$ que $\underline{L}$ \deleted{provient} est isomorphe à l'image inverse d'un Module birigidifié $\underline{M}$ sur $A\mathbin{\mathrm{x}}Q$, et pour la première assertion de \deleted{1}\handwritten{4}.7, on est donc ramené au cas $A=P$. Or $\underline L$ définit un morphisme de schémas ponctués}
\[
 \deleted{f\colon Q\handwritten{\longrightarrow}\underline{\mathrm{P}\deleted{\unsure{i|j}}\mathrm{ic}}_{P/k}\ ,}
\]
\deleted{qui se factorise nécessairement par le schéma abélien dual $P'$ de $P$. Soit $\underline{L}'$ l'image inverse du faisceau de Weil sur $P\mathbin{\mathrm{x}}P'$ par \deleted{$\mathrm{id}_{P}\mathbin{\mathrm{x}}\mathrm{id}$}}
\[
 \deleted{\deleted{\unsure{p|P}}\mathrm{id}_{P}\mathbin{\mathrm{x}}f\colon P\mathbin{\mathrm{x}}Q\handwritten{\longrightarrow}P\mathbin{\mathrm{x}}P'\ ,}
\]
\deleted{c'est un Module birigidifié qui, par définition de $f$, ne diffère de $\underline L$ que modulo l'image inverse d'un Module inversible $\underline M$ sur $P$. Comme $\underline L'$ et $\underline L$ ont une restriction triviale à $\deleted{\unsure{E|B}}P\mathbin{\mathrm{x}}e_{Q}$ , on en conclut que $\underline M$ est \deleted{\unsure{trivaal|trivaial}}trivial, donc que $\underline L$ est isomorphe à $\underline L'$. Or nous avons signalé dans VII 2.9.4 \deleted{\unsure{p|q}}que sur un produit de deux schémas abéliens, les correspondances divisorielles provenaient de biextensions. Il en résulte que $\underline L'$ , \deleted{\unsure{et en|et ea}} ou \deleted{\unsure{p|pe}}ce qui revient au même, $\underline L$ , provient d'une biextension.}

\deleted{Revenant au cas général ($P$ extension de $A$ par $L$), \deleted{\unsure{et|él}}il reste à prouver que (\deleted{1}\handwritten{4}.7.1) est bien un isomorphisme. Mais cela résulte de la suite exacte des $\mathrm{B}{\unsure{\bar{\mathrm{i}}|\mathrm{i}}}\mathrm{ext}^{i}$ ($i=0,1$) associée à la suite exacte}
\[
 \deleted{0\handwritten{\longrightarrow}L\handwritten{\longrightarrow}P\handwritten{\longrightarrow}A\handwritten{\longrightarrow}0\ ,}
\]
\deleted{et des relations}
\[
 \deleted{\operatorname{Biext}^{\unsure{o|0}}(L,Q;\underline G_{m})=0\ ,\quad
 \operatorname{Biext}^{1}(L\overset{\deleted{\mathrm{x}}}{,}Q;\underline G_{m})=0\ ,}
\]
\deleted{dont la première est un cas \deleted{\unsure{partioulier|partiwulier}}particulier de l'assertion de rigidité faite dans \deleted{1}\handwritten{4}.\handwritten{3} , et la deuxième n'est autre que le premier cas de \deleted{1}\handwritten{4}.6.}

\deleted{\uline{Corollaire} \deleted{1}\handwritten{4}.8. Sous les conditions de \deleted{1}\handwritten{4}.7, supposons que $Q$ soit extension d'un groupe algébrique $B$ (\unsure{nécesairement|nécessairement} lisse et connexe) par un}
% A short vertical pen stroke brackets the two lines of the corollary.
% archive sheet number: 49
\scanpage{51}
\hfill SGA 5
% archive sheet number: 50
\scanpage{57}
% letterhead
\begin{center}
INSTITUT\\
DES HAUTES ÉTUDES SCIENTIFIQUES\\[-.4ex]
\rule{1.4cm}{.4pt}
\end{center}

SGA 6 \uline{$\overline{\mathrm{V}\deleted{\unsure{\mathrm{I}|\mathrm{II}}}}$}\qquad Donner interprétation de la relation entre ``\unsure{suites|suite} mult. de polynômes'' de \unsure{Hirzebruch|Hizebruch}, et séries formelles, par \unsure{la 'opération|la l'opération|l'opération} $f*g$ \ldots

\noindent\rule{\linewidth}{.4pt}

Introduction, page de garde

Mentionner Serre et \unsure{Jouanolou} sur page de garde ?
% archive sheet number: 56
\scanpage{58}
Pour Exp II, \inserted{2.2.6} signaler

SGA 2 \quad V 3.5.

p. 10.
\margin{\unsure{ONG}}
% Faint large outlined letters in the paper, at middle left: ONG (left edge cropped).
% archive sheet number: 57
\scanpage{59}
\unsure{Posons}
\[
 \bar i\,\bar j=\tau
\]

Pour $\forall\ E\in\operatorname{Ob}C$, \unsure{nous allons} définir une application
\[
 \tikz[baseline=(a.base)]{\node(a) at (0,0){$*$};\draw (0,0) ellipse (.65em and .36em);}\qquad
 \tau(E)\times E\xrightarrow{{\unsure{\varphi|\phi}}_{E}}E
\]
qui \unsure{fera|fer.} de $E$ un objet formellement principal sous l'objet-groupe $\tau(E)$. \deleted{Il \unsure{revient|revient à|prends}} \inserted{\unsure{Pour ceci|Pour voir}} \unsure{si|Si} \ldots

\[
 \deleted{\text{\unsure{on|ou}}\quad\unsure{\bar{\imath}|=}\quad\unsure{\cdot|i}}\quad
 \unsure{\cdot|\ldots}\longrightarrow{J}\longrightarrow E\xrightarrow{{\unsure{p|\rho|\varphi}}_{E}}B
\]
\unsure{on aura|on a} $\unsure{\tau(E)|\tau(F)}=D_{\unsure{B|E}}(\unsure{J|\mathcal{J}|I})$,
% The bottom edge cuts into this line, including the argument of tau and the subscript of D.
% archive sheet number: 58
\scanpage{60}
% letterhead
\begin{center}
INSTITUT\\
DES HAUTES ÉTUDES SCIENTIFIQUES\\[-.4ex]
\rule{1.4cm}{.4pt}
\end{center}

SGA 6 $\overline{\underline{\mathrm{VII}}}$ \qquad \unsure{manque|marquer} la formule de self-intersection pour $\mathrm{Gr}(X),\mathrm{Gr}(Y)$

(après \deleted{\unsure{2|4|6}}4.\deleted{\unsure{5|6}}\unsure{6|5})

Peut-être \unsure{fin|bien} de l'Exp ?
% archive sheet number: 59
\scanpage{61}
% letterhead
\begin{center}
INSTITUT\\
DES HAUTES ÉTUDES SCIENTIFIQUES\\[-.4ex]
\rule{1.4cm}{.4pt}
\end{center}

SGA 6

\unsure{Structure} additive explicite de l'anneau $\unsure{\mathrm{Gr}|G_{r}}(\widetilde X)\otimes_{\unsure{\mathbb Z|Z|2}}\mathbb Q$ d'un schéma \unsure{éclaté|éclati}\ldots
% archive sheet number: 60
\scanpage{62}
% letterhead
\begin{center}
INSTITUT\\
DES HAUTES ÉTUDES SCIENTIFIQUES\\[-.4ex]
\rule{1.4cm}{.4pt}
\end{center}

\uline{Exp}. X \qquad Tenir compte corrections \unsure{de|des} Illusie (faire relire par Berthelot)
% archive sheet number: 61
\scanpage{63}
% letterhead
\begin{center}
INSTITUT\\
DES HAUTES ÉTUDES SCIENTIFIQUES\\[-.4ex]
\rule{1.4cm}{.4pt}
\end{center}

SGA \deleted{\unsure{0|5}}6 $\overline{\underline{\mathrm{III}}}$ 1, p. 15, signaler en footnote \margin{\deleted{\unsure{15|14}}}qu'un morphisme plat loc. de prés. finie est \deleted{\unsure{plat|pd}}parfait, et à \unsure{fortiori|fortiori.} pseudo-cohérent.
% archive sheet number: 62
